\clearpage
\phantomsection% 
\addcontentsline{toc}{chapter}{ABSTRAK}
\thispagestyle{fancy}

\begin{center}
	\large \bfseries \MakeUppercase{Abstrak}\\
%	\normalsize \normalfont {\thetitle}\\
%	\normalsize \normalfont {\theauthor}\\
	\bigskip
	
	\normalsize \normalfont \justifying \singlespacing
	Perkembangan kecerdasan buatan generatif telah meningkatkan kemampuan pembuatan video \textit{deepfake} yang semakin realistis, sehingga menimbulkan tantangan dalam keamanan informasi dan kepercayaan digital. Deteksi \textit{deepfake} pada data video dunia nyata menjadi semakin sulit karena adanya variasi kualitas visual, pose, pencahayaan, resolusi, dan karakteristik temporal. Penelitian ini menganalisis kontribusi pendekatan spasiotemporal berbasis VideoMAE dibandingkan pendekatan spasial berbasis Vision Transformer (ViT) dalam klasifikasi video \textit{deepfake}. Selain itu, penelitian ini menganalisis sensitivitas kedua pendekatan terhadap perubahan konfigurasi pelatihan dan karakteristik input, serta mengevaluasi konsistensi efektivitasnya pada skenario lintas-\textit{dataset}. Dataset utama yang digunakan adalah \textit{WildDeepfake}, sedangkan \textit{Celeb-DF v2} digunakan sebagai dataset eksternal untuk evaluasi \textit{zero-shot}. Tahapan penelitian meliputi \textit{preprocessing} data, pembentukan klip 46 \textit{frame}, \textit{strided temporal sampling}, pelatihan model, pengujian konfigurasi, dan evaluasi menggunakan \textit{accuracy}, \textit{precision}, \textit{recall}, \textit{F1-score}, AUC-ROC, dan \textit{confusion matrix}. Pada eksperimen \textit{baseline} di \textit{WildDeepfake}, model spasiotemporal memperoleh \textit{accuracy} sebesar 85,8\%, \textit{F1-score} sebesar 84,5\%, dan \textit{AUC} sebesar 92,3\%, sedangkan model spasial memperoleh \textit{accuracy} sebesar 79,2\%, \textit{F1-score} sebesar 76,7\%, dan \textit{AUC} sebesar 86,8\%. Hasil tersebut menunjukkan bahwa informasi antar-\textit{frame} memberikan kontribusi positif pada evaluasi internal. Namun, pada evaluasi lintas-\textit{dataset} terhadap \textit{Celeb-DF v2}, model spasial memperoleh \textit{accuracy} sebesar 72,9\% dan \textit{AUC} sebesar 77,7\%, sedangkan model spasiotemporal memperoleh \textit{accuracy} sebesar 64,2\% dan \textit{AUC} sebesar 72,6\%. Dengan demikian, pendekatan spasiotemporal lebih efektif pada evaluasi internal \textit{WildDeepfake}, sedangkan pendekatan spasial menunjukkan generalisasi yang lebih baik pada skenario \textit{zero-shot} lintas-\textit{dataset}. \par

    \vspace{20pt}
	\raggedright \textbf{Kata Kunci}: \textit{deepfake}, ViT, VideoMAE, spasiotemporal, \textit{WildDeepfake}
	
	\vfill
	
\end{center}
\clearpage
